% ID: SEL-D0072-OB-E03
% ESTADO: cotejada. Segunda pasada visual y borrador compilado revisados.
% Fuente: Selectividad/sel_2009_matematicas_aplicadas/mat apli CC SS06.pdf
% SHA-256: b371bdd6ba8b176301371f25945c7fb030327ff4501a6a02aa6f109a3a4f53e1
% Página: 2; opción B; ejercicio original 3.

\noindent\underline{Parte I}

El $70\%$ de los visitantes de un museo son españoles. El $49\%$ son españoles y mayores de edad. De los que no son españoles, el $40\%$ son menores de edad.
\begin{enumerate}[label=\alph*)]
\item \textbf{(1 punto)} Si se escoge, al azar, un visitante de este museo, \textquestiondown{}cuál es la probabilidad de que sea mayor de edad?
\item \textbf{(1 punto)} Se ha elegido, aleatoriamente, un visitante de este museo y resulta que es menor de edad, \textquestiondown{}cuál es la probabilidad de que no sea español?
\end{enumerate}
\par\noindent\begin{minipage}{\linewidth}
\noindent\underline{Parte II}

Los jóvenes andaluces duermen un número de horas diarias que se distribuye según una ley Normal de media desconocida, $\mu$, y desviación típica 2 horas. A partir de una muestra de 64 jóvenes se ha obtenido una media de 7 horas.
\begin{enumerate}[label=\alph*)]
\item \textbf{(1 punto)} Halle un intervalo de confianza, al $97\%$, para la media poblacional $\mu$.
\item \textbf{(1 punto)} Manteniendo la misma confianza, \textquestiondown{}cuál deberá ser el tamaño mínimo de la muestra para estimar la media de horas de sueño, cometiendo un error máximo de $0.25$ horas?
\end{enumerate}
\end{minipage}\par
